\documentclass[11pt]{article}
\usepackage[letterpaper,margin=1in]{geometry}
\usepackage{amsmath,amsthm,mathtools}
\usepackage{newtxtext,newtxmath}
\usepackage{booktabs,array}
\usepackage{enumitem}
\usepackage[hidelinks]{hyperref}

\newtheorem{theorem}{Theorem}
\newtheorem{lemma}{Lemma}
\renewcommand{\thelemma}{S\arabic{lemma}}
\renewcommand{\theequation}{S\arabic{equation}}
\renewcommand{\thesection}{S\arabic{section}}
\renewcommand{\thetable}{S\arabic{table}}
\renewcommand{\thesubsection}{\thesection.\arabic{subsection}}

\newcommand{\NE}{N^{\mathrm{E}}}
\newcommand{\NC}{N^{\mathrm{C}}}
\newcommand{\calF}{\mathcal{F}}
\newcommand{\calB}{\mathcal{B}}
\newcommand{\calS}{\mathcal{S}}
\newcommand{\calI}{\mathcal{I}}
\newcommand{\calR}{\mathcal{R}}
\newcommand{\calP}{\mathcal{P}}
\newcommand{\Z}{\mathbb{Z}}
\newcommand{\Cs}{C^{\star}}
\newcommand{\CLP}{C^{\mathrm{LP}}}
\newcommand{\dmin}{\vphantom{\underline{d}}\smash[b]{\underline{d}}}
\newcommand{\dmax}{\overline{d}}
\newcommand{\etamin}{\vphantom{\underline{\eta}}\smash[b]{\underline{\eta}}}
\newcommand{\conv}{\operatorname{conv}}
\newcommand{\disc}{\operatorname{disc}}

\title{Supplementary Material for\\ ``Cast-Break-Free Demand Response of Electric Arc Furnace Fleets:\\ Exact Counting Polytopes and Certified Rolling Dispatch''}
\author{Kaidi Huang}
\date{}

\begin{document}
\maketitle

This document contains the proofs of Theorems~1 to~4 of the paper, the rows of the counting polytope for the live
state and the end of the horizon, and an example for Theorem~1(ii). Section~\ref{sec:model} restates the model and the
notation, so that the document can be read on its own. Equations and lemmas numbered with the prefix S belong to this
document, and Theorems~1 to~4 are numbered as in the paper.

\section{Model and Notation}\label{sec:model}

\subsection{Zero-Break Schedules of One Plant}

Where a single plant is concerned, we drop the plant index. Time is divided into slots $t$. A plant has $\NE$ identical
furnaces and $\NC$ identical casters, and any of its ladles may be cast on any caster. A heat started at slot $s$
occupies a furnace for $d(s)$ slots, taps its ladle at slot $s+d(s)$, and frees the furnace $g$ slots after the tap. A
cast started at slot $z$ occupies a caster for $c(z)$ slots. Throughout, $t+d(t)$ and $t+c(t)$ are non-decreasing in
$t$, so a later start never ends earlier.

A schedule gives every heat $h$ a start $s_h$ and a casting start $z_h$ for its ladle, and it chooses the pairs
$\calP^{\mathrm E}$ of consecutive heats on a furnace and the pairs $\calP^{\mathrm C}$ of consecutive ladles on a
caster. The schedule is zero-break if
\begin{subequations}\label{eq:zb}
\begin{align}
& s_{h'}\ge s_h+d(s_h)+g, && (h,h')\in\calP^{\mathrm E}, \label{eq:zb-furn}\\
& z_h+c(z_h)\le z_{h'}\le z_h+c(z_h)+b, && (h,h')\in\calP^{\mathrm C}, \label{eq:zb-cast}\\
& s_h+d(s_h)\le z_h\le s_h+d(s_h)+w, && \text{for every heat } h. \label{eq:zb-ladle}
\end{align}
\end{subequations}
Thus a furnace is free $g$ slots after each tap, a caster stays busy during a cast and idles for at most $b$ slots
before the next one, and every ladle is cast within $w$ slots of its tap. These are the zero-break constraints of the
fleet model in Section~II of the paper.

\subsection{Counters, Live State and the Counting Polytope}

Let $x(t)$ and $y(t)$ count the heats and the casts started at or before slot $t$, and let
\begin{equation}\label{eq:theta}
\theta(t)=\max\{s:\ s+d(s)\le t\},\qquad r(t)=x(t)-x\big(\theta(t)\big).
\end{equation}
Since $t+d(t)$ is non-decreasing, heats tap in the order of their starts, $\theta(t)$ is the last start whose heat has
tapped by slot $t$, and $r(t)$ is the number of heats in progress at slot $t$.

The notice is given at $t=0$. The live state enters as one last start per furnace and per caster, placed so that its
heat or cast ends at the actual tap or cast end. Heats are indexed $k=1,2,\dots$ by start from these last starts on, so
heats $1$ to $n$ are the last heats of the $n$ furnaces that are idle at the notice and have tapped before it, and the
heats $k>n$ tap from $t=0$ on. Casts are indexed likewise, so casts $1$ to $\NC$ are the last casts of the casters.
With $Q$ ladles waiting at the notice, first-in first-out casting gives the waiting ladles to casts $\NC+1$ to $\NC+Q$
and the ladle of heat $k>n$ to cast $k+q$, where
\begin{equation}\label{eq:q}
q=Q+\NC-n.
\end{equation}

The counting polytope $\calF$ collects the counters that satisfy, at every slot $t$ from the last starts of the live
state to the horizon $H$,
\begin{subequations}\label{eq:F}
\begin{align}
& x(t)\ge x(t-1),\qquad y(t)\ge y(t-1), \label{eq:F-mono}\\
& x\big(t+d(t)+g-1\big)-x(t-1)\le\NE, \label{eq:F-reuse}\\
& y\big(t+c(t)-1\big)-y(t-1)\le\NC, \label{eq:F-busy}\\
& y(t)-x\big(\theta(t)\big)\le q, \label{eq:F-prec}\\
& y\big(t+d(t)+w\big)-x(t)\ge q, \label{eq:F-wait}\\
& y\big(t+c(t)+b\big)-y(t)\ge\NC, \label{eq:F-nobreak}
\end{align}
\end{subequations}
together with the following boundary rows.
\begin{itemize}[leftmargin=1.5em]
\item The live state fixes the counters before $t=0$. Row \eqref{eq:F-prec} is imposed for $t\ge0$, and row
\eqref{eq:F-wait} for the heats tapped from $t=0$ on.
\item The $j$-th oldest waiting ladle, tapped at $\tau_j\le0$, adds the row $y(\tau_j+w)\ge\NC+j$.
\item Beyond $H$, rows \eqref{eq:F-reuse} and \eqref{eq:F-busy} cap their arguments at $H$, that is, they read
$x(\min\{t+d(t)+g-1,H\})$ and $y(\min\{t+c(t)-1,H\})$, and the other rows are dropped.
\end{itemize}
Every row bounds the difference of two counters, or a single counter, by an integer, that is, every row is a
difference constraint. The round-robin assignment puts heat $k$ on furnace $k\bmod\NE$, cast $k$ on caster $k\bmod\NC$,
and the ladle of heat $k$ in cast $k+q$.

A call window consists of $m$ blocks $\beta\in\calB$ of three slots. The block load of a plant whose furnaces have power
$P$ is its average power in the block,
\begin{equation}\label{eq:load}
\ell_\beta=\frac P3\sum_{t\in\beta}r(t).
\end{equation}

\section{Proof of Theorem 1}\label{sec:thm1}

\begin{theorem}[Exact Counting Polytope]
(i) The integer points of $\calF$ are exactly the counters of the plant's zero-break schedules, each realized by the
round-robin assignment. (ii) $\calF$ is an integral polytope, and the block loads \eqref{eq:load} map it onto the convex
hull of the block loads of these schedules.
\end{theorem}

We first translate the rows of $\calF$ into inequalities between the starts of individual heats and casts. For
integer, non-decreasing counters, let
\begin{equation}\label{eq:order}
s_k=\min\{t:\ x(t)\ge k\},\qquad v_k=s_k+d(s_k),\qquad z_k=\min\{t:\ y(t)\ge k\}
\end{equation}
be the start and the tap of heat $k$ and the start of cast $k$. Then $x(t)\ge k$ if and only if $s_k\le t$, and
$y(t)\ge k$ if and only if $z_k\le t$. Since $t+d(t)$ is non-decreasing, heat $k$ has tapped by slot $t$, that is,
$v_k\le t$, if and only if $s_k\le\theta(t)$, that is, if and only if $x(\theta(t))\ge k$.

\begin{lemma}\label{lem:order-stat}
For integer counters that satisfy \eqref{eq:F-mono}, rows \eqref{eq:F-reuse} to \eqref{eq:F-nobreak} hold at every slot
if and only if, for every $k$,
\begin{subequations}\label{eq:S}
\begin{align}
& s_{k+\NE}\ge v_k+g, \label{eq:S-furn}\\
& z_k+c(z_k)\le z_{k+\NC}\le z_k+c(z_k)+b, \label{eq:S-cast}\\
& v_k\le z_{k+q}\le v_k+w, \label{eq:S-ladle}
\end{align}
\end{subequations}
where each inequality is required when its starts lie in the horizon, and \eqref{eq:S-ladle} is required for $k>n$.
\end{lemma}

\begin{proof}
We treat the rows in three groups.

\emph{Row \eqref{eq:F-reuse} and inequality \eqref{eq:S-furn}.} Suppose \eqref{eq:S-furn} holds and fix $t$. Let
$k=x(t-1)+1$ be the first heat started at or after $t$, so $s_k\ge t$. By \eqref{eq:S-furn} and the monotonicity of
$t+d(t)$, heat $k+\NE$ starts at or after $v_k+g\ge t+d(t)+g$. Hence at most the $\NE$ heats $k,\dots,k+\NE-1$ start in
$[t,\,t+d(t)+g-1]$, which is \eqref{eq:F-reuse}. Conversely, row \eqref{eq:F-reuse} at $t=s_k$ gives
\[
x(v_k+g-1)\le x(s_k-1)+\NE\le k-1+\NE,
\]
so heat $k+\NE$ starts after $v_k+g-1$, which is \eqref{eq:S-furn}.

\emph{Rows \eqref{eq:F-busy}, \eqref{eq:F-nobreak} and inequality \eqref{eq:S-cast}.} The argument above, with $c$,
$0$ and the casts in place of $d$, $g$ and the heats, shows that \eqref{eq:F-busy} is equivalent to the left inequality
of \eqref{eq:S-cast}. Suppose the right inequality holds and fix $t$. With $k=y(t)$, we have $z_k\le t$, so cast
$k+\NC$ starts by $z_k+c(z_k)+b\le t+c(t)+b$. Hence $y(t+c(t)+b)\ge k+\NC$, which is \eqref{eq:F-nobreak}.
Conversely, row \eqref{eq:F-nobreak} at $t=z_k$ gives $y(z_k+c(z_k)+b)\ge y(z_k)+\NC\ge k+\NC$, which is the right
inequality of \eqref{eq:S-cast}.

\emph{Rows \eqref{eq:F-prec}, \eqref{eq:F-wait} and inequality \eqref{eq:S-ladle}.} Suppose the left inequality of
\eqref{eq:S-ladle} holds and fix $t\ge0$. Let $j=y(t)$. If $j\le\NC+Q$, then $y(t)-x(\theta(t))\le\NC+Q-n=q$, since
$x(\theta(t))\ge n$ for $t\ge0$. Otherwise heat $k=j-q>n$ feeds cast $j$, which starts by $t$, so $v_k\le z_j\le t$ and
$x(\theta(t))\ge k=y(t)-q$. Either way \eqref{eq:F-prec} holds. Conversely, row \eqref{eq:F-prec} at $t=z_{k+q}$ gives
$x(\theta(z_{k+q}))\ge y(z_{k+q})-q\ge k$, so heat $k$ taps by $z_{k+q}$. Suppose next that the right inequality of
\eqref{eq:S-ladle} holds, and fix a slot $t$ with $k=x(t)>n$. Then $s_k\le t$ and $v_k\le t+d(t)$, so cast $k+q$
starts by $v_k+w\le t+d(t)+w$, which gives $y(t+d(t)+w)\ge k+q$, that is, \eqref{eq:F-wait}. Conversely, row
\eqref{eq:F-wait} at $t=s_k$ gives $y(v_k+w)\ge x(s_k)+q\ge k+q$, that is, $z_{k+q}\le v_k+w$.
\end{proof}

\begin{proof}[Proof of Theorem 1]
\emph{(i), zero-break schedules satisfy the rows of $\calF$.} Take a zero-break schedule and its counters. Row
\eqref{eq:F-mono} holds because the counters count starts. Fix a slot $t$.
\begin{itemize}[leftmargin=1.5em]
\item If a furnace starts a heat at $s\in[t,\,t+d(t)+g-1]$, then \eqref{eq:zb-furn} and $s+d(s)\ge t+d(t)$ place its
next start at or after $t+d(t)+g$. Each furnace thus starts at most one heat in this window, which gives
\eqref{eq:F-reuse}.
\item In the same way, \eqref{eq:zb-cast} lets each caster start at most one cast in $[t,\,t+c(t)-1]$, which gives
\eqref{eq:F-busy}. It also makes each caster start at least one cast in $(t,\,t+c(t)+b]$, namely the successor of its
last cast started by $t$, which gives \eqref{eq:F-nobreak}.
\item The $y(t)-\NC$ casts started by $t$ after the last casts take distinct ladles. These ladles are among the $Q$
waiting ones and the $x(\theta(t))-n$ ladles tapped from $t=0$ on by slot $t$, so $y(t)-\NC\le Q+x(\theta(t))-n$, which
is \eqref{eq:F-prec} by \eqref{eq:q}.
\item By \eqref{eq:zb-ladle}, the waiting ladles and the ladles of the $x(t)-n$ heats $k>n$ started by $t$ are cast by
$t+d(t)+w$, so $y(t+d(t)+w)-\NC\ge Q+x(t)-n$, which is \eqref{eq:F-wait}. The rows of the waiting ladles hold for the
same reason.
\end{itemize}
Finally, the counters do not decrease after $H$, so capping the arguments at $H$ only weakens rows \eqref{eq:F-reuse}
and \eqref{eq:F-busy}.

\emph{(i), integer points of $\calF$ are zero-break schedules.} Let $(x,y)$ be an integer point of $\calF$, and let
\eqref{eq:order} give its starts. By Lemma~\ref{lem:order-stat}, they satisfy \eqref{eq:S}. Under the round-robin
assignment, heat $k+\NE$ is the next heat on the furnace of heat $k$, cast $k+\NC$ is the next cast on the caster of
cast $k$, and cast $k+q$ receives the ladle of heat $k$. Inequalities \eqref{eq:S-furn}, \eqref{eq:S-cast} and
\eqref{eq:S-ladle} are therefore \eqref{eq:zb-furn}, \eqref{eq:zb-cast} and \eqref{eq:zb-ladle} for the pairs and the
ladles of this assignment, and the rows of the waiting ladles give \eqref{eq:zb-ladle} for the ladles waiting at the
notice. The schedule is thus zero-break, and its counters are $(x,y)$.

\emph{(ii).} Each row of $\calF$ has at most one coefficient $+1$ and one coefficient $-1$ and an integer right side.
The constraint matrix is therefore the transpose of the node-arc incidence matrix of a directed graph, with additional
unit rows for the single-counter rows, and it is totally unimodular \cite{durr2011algorithmica}. Hence $\calF$ is an
integral polyhedron. It is bounded, because the counters start from the values fixed by the live state, do not
decrease by \eqref{eq:F-mono}, and grow by at most $\NE$ and $\NC$ per slot by \eqref{eq:F-reuse} and
\eqref{eq:F-busy}. So $\calF=\conv(\calF\cap\Z)$. The block loads \eqref{eq:load} are linear in the counters, so they
map $\conv(\calF\cap\Z)$ onto the convex hull of the block loads of $\calF\cap\Z$, which are the block loads of the
zero-break schedules by (i).
\end{proof}

\section{An Example for Theorem 1(ii)}\label{sec:example}

Every edge of a generalized polymatroid (g-polymatroid) is parallel to a unit vector $e_j$ or to a difference
$e_j-e_k$, and so is every edge of its projection onto a subset of coordinates \cite{mukhi2026tsg}. The following plant
shows that the image in Theorem~1(ii) need not have this property.

Consider one furnace and one caster with $d=2$, $g=0$, $c=3$ and $b=w=1$. At the notice, the furnace is idle, no ladle
is waiting, and the cast in progress ends at $t=1$. The window starts at the notice, $t_0=0$, and $H\ge10$. The first
cast must then start at slot $2$, with the ladle of a heat started at slot $0$, and every later cast starts three or
four slots after the previous one, with the ladle of a heat tapped at its start or one slot earlier, while heats on the
furnace do not overlap. Enumerating these schedules gives, in units of $P/3$, the nine load vectors of the first three
blocks listed in Table~\ref{tab:example}.

\begin{table}[h]
\centering
\caption{Load vectors $(\ell_1,\ell_2,\ell_3)$ of the first three blocks, in units of $P/3$, and the value of
$\ell_1+2\ell_2+\ell_3$.}
\label{tab:example}
\begin{tabular}{ccccccccc}
\toprule
$(2,2,1)$ & $(3,1,2)$ & $(2,2,2)$ & $(3,1,3)$ & $(3,2,1)$ & $(2,2,3)$ & $(2,3,1)$ & $(3,2,2)$ & $(2,3,2)$\\
$7$ & $7$ & $8$ & $8$ & $8$ & $9$ & $9$ & $9$ & $10$\\
\bottomrule
\end{tabular}
\end{table}

For example, heats started at slots $0$, $3$ and $8$, cast at slots $2$, $6$ and $10$, give $(2,2,1)$, and heats started
at slots $0$, $2$, $6$ and $9$, cast at slots $2$, $5$, $8$ and $11$, give $(3,1,2)$. The linear function
$\ell_1+2\ell_2+\ell_3$ equals $7$ at these two vectors and at least $8$ at the other seven, so its minimum over the
projected hull is attained exactly on the segment between them, which is therefore an edge. The direction of this edge,
$(1,-1,1)$, is parallel to no $e_j$ and to no $e_j-e_k$. The image of $\calF$ is therefore not a g-polymatroid.

\section{Proof of Theorem 2}\label{sec:thm2}

\subsection{The Arc LP}

A heat needs energy $E$, counted in MW-slots, and draws any power between $\underline P$ and $\overline P$ in each slot
that it occupies the furnace. A heat started at $s$ and tapped at $v$ is an arc $(s,v)$, admissible if
$(v-s)\underline P\le E\le(v-s)\overline P$, so the admissible durations $v-s$ form an interval. At slot prices $\pi_t$,
its cheapest power profile costs
\begin{equation}\label{eq:kappa}
\kappa(s,v)=\min_p\Big\{\sum_{t=s}^{v-1}\pi_tp_t:\ \underline P\le p_t\le\overline P,\ \sum_{t=s}^{v-1}p_t=E\Big\}.
\end{equation}
The arc LP has one variable $\xi(s,v)\ge0$ per admissible arc, counts the starts and the taps by
\begin{equation}\label{eq:arc}
x(t)=\sum_{s\le t}\sum_v\xi(s,v),\qquad \hat x(t)=\sum_{v\le t}\sum_s\xi(s,v),
\end{equation}
and imposes, with the tap now a decision, the rows
\begin{equation}\label{eq:arcrows}
x(t)-\hat x(t-g)\le\NE,\qquad y(t)-\hat x(t)\le q,\qquad y(t+w)-\hat x(t)\ge q,
\end{equation}
together with the cast rows \eqref{eq:F-mono}, \eqref{eq:F-busy} and \eqref{eq:F-nobreak} and the boundary rows. It
minimizes $\sum_{(s,v)}\kappa(s,v)\xi(s,v)$ plus any cost that is linear in $y$.

\begin{theorem}[Exact Power Band]
For all prices $\pi$, the arc LP has an integral optimal solution, which is a cheapest zero-break schedule with power
within the band.
\end{theorem}

The proof shows that the arc costs are Monge, uses this to remove nested arcs from an optimal solution, and then rounds
the solution without loss.

\begin{proof}
\emph{Step 1: closed form of the arc costs.} Write $p_t=\underline P+(\overline P-\underline P)\omega_t$ with
$\omega_t\in[0,1]$. Then
\begin{equation}\label{eq:kappa-split}
\kappa(s,v)=\underline P\sum_{t=s}^{v-1}\pi_t+(\overline P-\underline P)\,\phi\big(\{s,\dots,v-1\},\nu_{v-s}\big),
\end{equation}
where
\[
\phi(I,\nu)=\min\Big\{\sum_{t\in I}\pi_t\omega_t:\ 0\le\omega_t\le1,\ \sum_{t\in I}\omega_t=\nu\Big\},\qquad
\nu_d=\frac{E-d\underline P}{\overline P-\underline P}.
\]
The minimum in $\phi$ fills the cheapest slots of $I$ first, so $\phi(I,\cdot)$ is piecewise linear on $[0,|I|]$ with
slopes equal to the sorted prices of $I$, and hence convex. Moreover $\nu_{d+1}=\nu_d-\delta$ with
$\delta=\underline P/(\overline P-\underline P)\ge0$. The first term of \eqref{eq:kappa-split} is a difference of
cumulative prices at $v$ and $s$, so it is modular in $(s,v)$ and cancels in the inequalities below.

\emph{Step 2: a local Monge inequality.} Fix $s<v$ such that the arcs $(s,v+1)$ and $(s+1,v)$ are admissible. Then
the arcs $(s,v)$ and $(s+1,v+1)$, whose duration lies in between, are admissible as well. Let $I=\{s+1,\dots,v-1\}$ and
$\nu=\nu_{v-s}$. The arcs $(s,v)$ and $(s+1,v+1)$ use the slots $I\cup\{s\}$ and $I\cup\{v\}$ with the parameter $\nu$,
the arc $(s,v+1)$ uses $I\cup\{s,v\}$ with $\nu-\delta$, and the arc $(s+1,v)$ uses $I$ with $\nu+\delta$. Let
$\alpha$ and $\alpha'$ be the weights of slots $s$ and $v$ in an optimal solution of $\phi(I\cup\{s,v\},\nu-\delta)$,
so that
\begin{equation}\label{eq:m1}
\phi(I\cup\{s,v\},\nu-\delta)=\pi_s\alpha+\pi_v\alpha'+\phi(I,\nu-\delta-\alpha-\alpha').
\end{equation}
The four numbers $\nu-\alpha$, $\nu-\alpha'$, $\nu-\delta-\alpha-\alpha'$ and $\nu+\delta$ lie in $[0,|I|]$. The smallest,
$\nu-\delta-\alpha-\alpha'$, is nonnegative because $\alpha+\alpha'\le\nu-\delta$ in \eqref{eq:m1}, and the largest satisfies
$\nu+\delta=\nu_{v-s-1}\le v-s-1=|I|$ because $(s+1,v)$ is admissible. Putting the weight $\alpha$ on slot $s$ and $\alpha'$
on slot $v$ therefore gives feasible points of the two problems on the left of

\begin{equation}\label{eq:m2}
\phi(I\cup\{s\},\nu)+\phi(I\cup\{v\},\nu)\le\pi_s\alpha+\pi_v\alpha'+\phi(I,\nu-\alpha)+\phi(I,\nu-\alpha').
\end{equation}
The pairs $(\nu-\alpha,\nu-\alpha')$ and $(\nu-\delta-\alpha-\alpha',\nu+\delta)$ have the same sum, and the second
pair contains the first between its two members, so convexity of $\phi(I,\cdot)$ gives
\begin{equation}\label{eq:m3}
\phi(I,\nu-\alpha)+\phi(I,\nu-\alpha')\le\phi(I,\nu-\delta-\alpha-\alpha')+\phi(I,\nu+\delta).
\end{equation}
Combining \eqref{eq:m1}, \eqref{eq:m2} and \eqref{eq:m3} yields
$\phi(I\cup\{s\},\nu)+\phi(I\cup\{v\},\nu)\le\phi(I\cup\{s,v\},\nu-\delta)+\phi(I,\nu+\delta)$, which by
\eqref{eq:kappa-split} reads
\begin{equation}\label{eq:local-monge}
\kappa(s,v)+\kappa(s+1,v+1)\le\kappa(s,v+1)+\kappa(s+1,v).
\end{equation}

\emph{Step 3: the Monge inequality for nested arcs.} Let $s_1<s_2$ and $v_1<v_2$, and let the nested arcs $(s_1,v_2)$
and $(s_2,v_1)$ be admissible. Every arc $(s,v)$ with $s_1\le s\le s_2$ and $v_1\le v\le v_2$ has a duration between
$v_1-s_2$ and $v_2-s_1$ and is therefore admissible. Summing \eqref{eq:local-monge} over $s=s_1,\dots,s_2-1$ and
$v=v_1,\dots,v_2-1$, the inner terms cancel, and
\begin{equation}\label{eq:monge}
\kappa(s_1,v_1)+\kappa(s_2,v_2)\le\kappa(s_1,v_2)+\kappa(s_2,v_1).
\end{equation}

\emph{Step 4: an optimal solution without nested arcs.} The heats in progress at the notice have fixed arcs, which only
shift the integer right sides of the rows, so we consider the arcs of the new heats alone. Let $\xi$ be optimal, and
suppose that it puts positive mass on nested arcs $(s_1,v_2)$ and $(s_2,v_1)$. Moving the mass
$\varepsilon=\min\{\xi(s_1,v_2),\xi(s_2,v_1)\}$ from these arcs to $(s_1,v_1)$ and $(s_2,v_2)$ leaves the start and tap
counters $x$ and $\hat x$ unchanged, so every row still holds. By \eqref{eq:monge}, the move does not raise the cost,
and it lowers $\sum_{(s,v)}\xi(s,v)(v-s)^2$ by $2\varepsilon(s_2-s_1)(v_2-v_1)>0$. An optimal solution that minimizes
$\sum_{(s,v)}\xi(s,v)(v-s)^2$ among all optimal solutions therefore puts no mass on nested arcs.

\emph{Step 5: lossless rounding.} Let $\xi$ be such a solution. Its support has no nested pair, so its arcs can be
ordered as $a_1,\dots,a_R$ with non-decreasing starts and non-decreasing taps. With $M_j=\sum_{i\le j}\xi(a_i)$, the
counters \eqref{eq:arc} take the values $x(t)=M_{J(t)}$ and $\hat x(t)=M_{\hat J(t)}$, where $J(t)$ and $\hat J(t)$ are
the largest indices of the arcs that start and that tap by slot $t$. All rows of the arc LP are difference constraints
in $x$, $\hat x$ and $y$ with integer right sides, and for every $u\in[0,1)$ and every such row $a-a'\le h$,
\begin{equation}\label{eq:floor}
\lfloor a+u\rfloor-\lfloor a'+u\rfloor\le\lceil a-a'\rceil\le h,
\end{equation}
because the left side is an integer smaller than $a-a'+1$. The shifted floors $\lfloor x(t)+u\rfloor$, $\lfloor\hat
x(t)+u\rfloor$ and $\lfloor y(t)+u\rfloor$ thus keep every row \cite{bansal2008sicomp}. The $k$-th start of the
rounded counters is the start of the arc $a_{j_k}$ with $j_k=\min\{j:\ M_j\ge k-u\}$, and the $k$-th tap is the tap of
the same arc, since starts and taps accumulate the masses $M_j$ in the same order. Each rounded heat is therefore an
arc of the support, run at its cheapest power profile, and the rounded counters form a zero-break schedule with power
within the band by the argument of Theorem~1. For $u$ uniform on $[0,1)$, the expected number of rounded heats on arc
$a_j$ is $\xi(a_j)$, and $\mathbb E\lfloor y(t)+u\rfloor=y(t)$, so the expected cost of the rounded schedule equals the
LP value. Some $u$ therefore gives an integral schedule whose cost is at most the LP value, and this schedule is an
integral optimal solution.
\end{proof}

\section{Proof of Theorem 3}\label{sec:thm3}

\subsection{Fleet LP, Shifted Floors and Certificates}

A fleet $\calI$ of plants answers a call of $m$ blocks $\beta\in\calB$, against the baseline $\bar L_\beta$ of each
block. Plant $i$ has its own data, counting polytope $\calF_i$ and furnace power $P_i$, and its block loads
$\ell_{i,\beta}$ follow \eqref{eq:load}. The counting integer program is
\begin{equation}\label{eq:ip}
\Cs=\min\Big\{C:\ (x_i,y_i)\in\calF_i\cap\Z,\ \sum_{i\in\calI}\ell_{i,\beta}-\bar L_\beta\le C,\ \beta\in\calB,\
i\in\calI\Big\},
\end{equation}
and its LP relaxation, the fleet LP, has optimal value $\CLP\le\Cs$. For an optimal solution of the fleet LP with block
loads $\ell^{\mathrm{LP}}_{i,\beta}$, let $\CLP_\beta=\sum_i\ell^{\mathrm{LP}}_{i,\beta}-\bar L_\beta$ be the LP excess
of block $\beta$, so $\CLP_\beta\le\CLP$, and let $\calS$ be the set of plants with a fractional counter.

The shifted floors of plant $i$ are
\begin{equation}\label{eq:shift}
x_i^u(t)=\lfloor x_i(t)+u_i\rfloor,\qquad y_i^u(t)=\lfloor y_i(t)+u_i\rfloor,
\end{equation}
with one shift $u_i\in[0,1)$ per plant, and $\ell^u_{i,\beta}$ denotes the block load of the rounded plant. For
independent shifts uniform on $[0,1)$, let
\begin{equation}\label{eq:U}
U(\Delta,\lambda)=\sum_{\beta\in\calB}e^{-\lambda(\Delta+\CLP-\CLP_\beta)}\prod_{i\in\calS}
\mathbb E\,e^{\lambda(\ell^u_{i,\beta}-\ell^{\mathrm{LP}}_{i,\beta})}.
\end{equation}
A margin $\Delta$ is a certificate if $U(\Delta,\lambda)<1$ for some $\lambda>0$. Given a certificate and its
$\lambda$, the derandomized rounding fixes the shifts of the plants in $\calS$ one at a time, each to minimize $U$ with
the earlier shifts fixed and the later ones random. Let $P_{\max}=\max_iP_i$.

\begin{theorem}[Fleet Rounding]
(i) For every certificate $\Delta$, the derandomized rounding finds zero-break schedules with a peak below
$\CLP+\Delta$ in polynomial time. At an optimal vertex of the fleet LP, $|\calS|\le m-1$, and every
$\Delta>\big(\tfrac12(m-1)\ln m\big)^{1/2}P_{\max}$ is a certificate. (ii) For single-furnace plants with
start-dependent heat and casting times and for arbitrary baselines, $\Cs-\CLP$ can reach $\Omega(\sqrt m\,P_{\max})$,
and for some constant $\epsilon>0$, no polynomial algorithm guarantees a peak within $\epsilon\sqrt m\,P_{\max}$ of $\Cs$
unless $\mathsf P=\mathsf{NP}$.
\end{theorem}

\subsection{Proof of Part (i)}

\begin{proof}[Proof of Theorem 3(i)]
\emph{Step 1: the shifted floors are zero-break schedules.} Every row of $\calF_i$ is a difference constraint, so by
\eqref{eq:floor} every shifted floor of a point of $\calF_i$ lies in $\calF_i\cap\Z$ \cite{bansal2008sicomp}, and by
Theorem~1(i) it is the counter of a zero-break schedule.

\emph{Step 2: the rounded loads are unbiased and bounded.} For $u$ uniform on $[0,1)$ and any real $a$, $\lfloor
a+u\rfloor$ equals $\lceil a\rceil$ with probability $a-\lfloor a\rfloor$ and $\lfloor a\rfloor$ otherwise, so
$\mathbb E\lfloor a+u\rfloor=a$. Every running count $r_i^u(t)=x_i^u(t)-x_i^u(\theta_i(t))$ therefore takes the value
$\lfloor r_i(t)\rfloor$ or $\lceil r_i(t)\rceil$ and has mean $r_i(t)$. By \eqref{eq:load}, $\mathbb
E\,\ell^u_{i,\beta}=\ell^{\mathrm{LP}}_{i,\beta}$, and $\ell^u_{i,\beta}$ stays in an interval of length
$\tfrac{P_i}{3}\cdot3=P_i$ as $u_i$ varies. A plant outside $\calS$ has integral counters and keeps its LP loads.

\emph{Step 3: the estimator bounds the probability of a high peak.} Fix a block $\beta$ and $\lambda>0$. Markov's
inequality applied to $e^{\lambda(\sum_i\ell^u_{i,\beta}-\bar L_\beta)}$ and the independence of the shifts give
\[
\Pr\Big\{\sum_{i\in\calI}\ell^u_{i,\beta}-\bar L_\beta\ge\CLP+\Delta\Big\}
\le e^{-\lambda(\Delta+\CLP-\CLP_\beta)}\prod_{i\in\calS}\mathbb E\,e^{\lambda(\ell^u_{i,\beta}-\ell^{\mathrm{LP}}_{i,\beta})},
\]
and the union bound over the blocks shows that $U(\Delta,\lambda)$ bounds the probability that the rounded peak
reaches $\CLP+\Delta$.

\emph{Step 4: derandomization.} Order the plants of $\calS$ and let $U_j$ be $U$ with the shifts of the first $j$ plants
fixed, that is, with their expectations in \eqref{eq:U} replaced by the values at the fixed shifts. Then $U_0=U$, and
$U_{j}$ is the average of $U_{j+1}$ over the shift of plant $j+1$, so a shift that minimizes $U_{j+1}$ gives
$U_{j+1}\le U_j$. Once all shifts are fixed,
\[
U_{|\calS|}=\sum_{\beta\in\calB}\exp\Big(\lambda\Big(\sum_{i\in\calI}\ell^u_{i,\beta}-\bar L_\beta-\CLP-\Delta\Big)\Big)
\le U(\Delta,\lambda)<1,
\]
so every term is below one, and every block of the rounded fleet lies below $\CLP+\Delta$.

\emph{Step 5: running time.} As a function of $u_i$, the shifted floors of plant $i$ change only where $u_i$ crosses one
minus the fractional part of a counter. The block loads involve the counters $x_i(t)$ and $x_i(\theta_i(t))$ of the
$3m$ slots of the window, so $\ell^u_{i,\beta}$ takes $O(m)$ values, each expectation in \eqref{eq:U} is an average of
$O(m)$ terms, and each step of the derandomization compares $O(m)$ candidate shifts. The rounding thus takes time
polynomial in $m$, $|\calI|$ and $H$.

\emph{Step 6: at most $m-1$ fractional plants.} Let $(C,x,y)$ be an optimal vertex of the fleet LP, and let $G_i$ be the
face of $\calF_i$ that contains $(x_i,y_i)$ in its relative interior. The face $G=\prod_iG_i\times\mathbb R$ contains
the vertex in its relative interior, and the vertex is the only point of $G$ near it that satisfies the tight coupling
rows with equality, since the other coupling rows have slack. Hence $\dim G$ is at most the number of tight coupling
rows, which is at most $m$. Since the vertices of $\calF_i$ are integral by Theorem~1(ii), every plant in $\calS$ has
$\dim G_i\ge1$, so
\[
1+|\calS|\le1+\sum_i\dim G_i=\dim G\le m,
\]
that is, $|\calS|\le m-1$, as in \cite{vujanic2016automatica}.

\emph{Step 7: a certificate independent of the fleet size.} By Step 2 and Hoeffding's lemma, $\mathbb
E\,e^{\lambda(\ell^u_{i,\beta}-\ell^{\mathrm{LP}}_{i,\beta})}\le e^{\lambda^2P_i^2/8}$. With $\CLP_\beta\le\CLP$,
\[
U(\Delta,\lambda)\le m\exp\Big(-\lambda\Delta+\frac{\lambda^2}{8}\sum_{i\in\calS}P_i^2\Big),
\]
and the choice $\lambda=4\Delta/\sum_{i\in\calS}P_i^2$ gives
\[
U(\Delta,\lambda)\le m\exp\Big(-\frac{2\Delta^2}{\sum_{i\in\calS}P_i^2}\Big),
\]
which is below one when $\Delta^2>\tfrac12\ln m\sum_{i\in\calS}P_i^2$. At an optimal vertex, Step 6 gives
$\sum_{i\in\calS}P_i^2\le(m-1)P_{\max}^2$, so every $\Delta>\big(\tfrac12(m-1)\ln m\big)^{1/2}P_{\max}$ is a
certificate. If $\calS$ is empty, every $\Delta>0$ is a certificate.
\end{proof}

\subsection{Proof of Part (ii)}

The proof reduces the discrepancy of set systems to the fleet problem. For a set system $T_1,\dots,T_M$ on the
elements $1,\dots,N$ and signs $\chi\in\{\pm1\}^N$, let $\sigma_k=\sum_{i\in T_k}\chi_i$. The discrepancy of the set
system is $\disc(T)=\min_\chi\max_k|\sigma_k|$.

\begin{proof}[Proof of Theorem 3(ii)]
\emph{Step 1: blocks.} Give set $k$ the two blocks $A_k$, of slots $6k$, $6k+1$, $6k+2$, and $B_k$, of slots $6k+3$,
$6k+4$, $6k+5$, counted from the start of the window, so that $m=2M$.

\emph{Step 2: plants.} Element $i$ becomes plant $i$, with one furnace of power $P=P_{\max}$, one caster, $g=w=0$ and
$b=1$. Its heat times depend on the start slot modulo six as in Table~\ref{tab:gadget}, so $s+d(s)$ is non-decreasing
and taps fall only on slots congruent to $3$ or $5$ modulo six. Since $w=0$, every cast starts at a tap. Let
$K_i=\{k:\ i\in T_k\}$. The casting times make a cast started at slot $6k+3$ or $6k+5$, with $k\in K_i$, last until slot
$6k'+3$ or $6k'+5$, respectively, where $k'$ is the next set of $K_i$, and the casting times at the other slots keep
$z+c(z)$ non-decreasing. A lead-in before the window, with its own heat and casting times, lets the first cast of plant
$i$ in its first set start at either slot $3$ or slot $5$ of that set.

\begin{table}[h]
\centering
\caption{Heat times of the plants of the reduction.}
\label{tab:gadget}
\begin{tabular}{lcccccc}
\toprule
Start slot $s$ modulo six & $0$ & $1$ & $2$ & $3$ & $4$ & $5$\\
Heat time $d(s)$ & $3$ & $4$ & $7$ & $6$ & $5$ & $4$\\
Tap $s+d(s)$ modulo six & $3$ & $5$ & $3$ & $3$ & $3$ & $3$\\
\bottomrule
\end{tabular}
\end{table}

\emph{Step 3: one sign per plant, and two dominant schedules.} A cast that ends at a slot congruent to $3$ must be
followed, since $b=1$, by a cast that starts at that slot or the next one, and only the former carries a tap. The same
holds for $5$. The offset, $3$ or $5$, of the casts of plant $i$ is therefore the same in all sets of $K_i$, and we write
$\chi_i=+1$ for offset $3$ and $\chi_i=-1$ for offset $5$. In every set $k\in K_i$, plant $i$ thus needs a heat that taps
at slot $6k+3$ if $\chi_i=+1$, or at slot $6k+5$ if $\chi_i=-1$. By Table~\ref{tab:gadget}, the tap at $6k+5$ requires
the start $6k+1$, so the heat runs in slots $6k+1$ to $6k+4$. The tap at $6k+3$ comes from the start $6k$, with the
heat in slots $6k$ to $6k+2$, or from a start in $6k-4,\dots,6k-1$, with a heat that covers these three slots and
earlier ones. The two minimal schedules add to the blocks of set $k$
\[
\chi_i=+1:\ P\text{ to }A_k\text{ and }0\text{ to }B_k,\qquad
\chi_i=-1:\ \tfrac23P\text{ to }A_k\text{ and }\tfrac23P\text{ to }B_k,
\]
that is, $P(5+\chi_i)/6$ to $A_k$ and $P(1-\chi_i)/3$ to $B_k$, and nothing to the blocks of the sets that do not contain
$i$. Every other schedule of plant $i$ runs these heats or longer ones, or extra heats, so it is dominated block by
block by the minimal schedule with the same sign.

\emph{Step 4: the gap equals a scaled discrepancy.} Fix $C_0$ and set the baselines
\[
\bar L_{A_k}=\sum_{i\in T_k}\frac{5P}{6}-C_0,\qquad \bar L_{B_k}=\sum_{i\in T_k}\frac{P}{3}-C_0.
\]
The minimal schedules then give the excesses $C_0+P\sigma_k/6$ in $A_k$ and $C_0-P\sigma_k/3$ in $B_k$. By
Theorem~1(ii), the fleet LP may use any convex combination of the schedules of a plant, so it may take
$\chi\in[-1,1]^N$, and it attains $C_0$ at $\chi=0$. It cannot do better, because one of $P\sigma_k/6$ and $-P\sigma_k/3$
is nonnegative for every $k$, and every combination of schedules is dominated by the same combination of minimal ones.
Hence $\CLP=C_0$, and since every integer schedule is dominated by one with $\chi\in\{\pm1\}^N$,
\begin{equation}\label{eq:gap}
\Cs-\CLP=\min_{\chi\in\{\pm1\}^N}\max_k\max\Big\{\frac P6\sigma_k,\,-\frac P3\sigma_k\Big\},\qquad
\frac P6\disc(T)\le\Cs-\CLP\le\frac P3\disc(T).
\end{equation}

\emph{Step 5: the gap can grow as $\sqrt m$.} Let $H$ be a Hadamard matrix of order $N$, a power of two, whose first row
has all entries $+1$, and let $T_k=\{i:\ H_{ki}=+1\}$ for $k=1,\dots,N$. For $\chi\in\{\pm1\}^N$, let $\Sigma=\sum_i\chi_i$.
Then $\sigma_1=\Sigma$ and $(H\chi)_k=2\sigma_k-\Sigma$. Since $H^\top H=NI$,
$\|H\chi\|^2=N\|\chi\|^2=N^2$, so some $k$ has $|2\sigma_k-\Sigma|\ge\sqrt N$ and thus $|\sigma_k|\ge(\sqrt
N-|\Sigma|)/2$. Hence
\[
\max_k|\sigma_k|\ge\max\Big\{|\Sigma|,\ \frac{\sqrt N-|\Sigma|}{2}\Big\}\ge\frac{\sqrt N}{3},
\]
so $\disc(T)\ge\sqrt N/3$. With $M=N$ and $m=2N$, \eqref{eq:gap} gives $\Cs-\CLP\ge P\sqrt N/18=\Omega(\sqrt
m\,P_{\max})$.

\emph{Step 6: hardness.} Charikar, Newman and Nikolov \cite{charikar2011soda} show that, for set systems with
$\Theta(N)$ sets on $N$ elements, it is NP-hard to tell discrepancy zero from discrepancy at least $c\sqrt N$, for some
constant $c>0$. The construction above has size polynomial in $N$, and $m=2M=\Theta(N)$. If $\disc(T)=0$, then
$\Cs=\CLP$, and if $\disc(T)\ge c\sqrt N$, then $\Cs\ge\CLP+cP\sqrt N/6$ by \eqref{eq:gap}. Choose $\epsilon>0$ with
$\epsilon\sqrt m<c\sqrt N/6$ for all such systems. A polynomial algorithm that guaranteed a peak within $\epsilon\sqrt
m\,P$ of $\Cs$ would then tell the two cases apart, by comparing its peak with $\CLP+cP\sqrt N/6$, where an LP gives $\CLP$
in polynomial time. No such algorithm exists unless $\mathsf P=\mathsf{NP}$.
\end{proof}

\section{Proof of Theorem 4}\label{sec:thm4}

\subsection{Tap Envelopes and the Certified Plant System}

Each heat taps between $s+\dmin(s)$ and $s+\dmax(s)$, where $s$ is its start. At slot $t$, the observations and these
bounds place the tap of a heat started at $s$ in a residual interval $[s+\dmin_t(s),\,s+\dmax_t(s)]$, which narrows as
the heat runs and shrinks to one slot at its tap. When several heats start at $s$, the interval is the smallest one
that contains theirs. The latest and earliest tap envelopes
\begin{equation}\label{eq:env}
\bar\eta_t(s)=\max_{s'\le s}\big(s'+\dmax_t(s')\big),\qquad \etamin_t(s)=\min_{s'\ge s}\big(s'+\dmin_t(s')\big)
\end{equation}
are non-decreasing in $s$, and as $t$ grows, $\bar\eta_t$ can only move earlier and $\etamin_t$ only later.

The certified system $\calR_t$ of a plant is its counting polytope on the horizon fixed at the notice, with the
counters before $t$ fixed to the executed starts, and with rows \eqref{eq:F-reuse}, \eqref{eq:F-prec} and
\eqref{eq:F-wait} replaced, at every slot $s$, by
\begin{subequations}\label{eq:R}
\begin{align}
& x\big(\bar\eta_t(s)+g-1\big)-x(s-1)\le\NE, \label{eq:R-reuse}\\
& y(s)-x\big(\bar\theta_t(s)\big)\le q,\qquad \bar\theta_t(s)=\max\{s':\ \bar\eta_t(s')\le s\}, \label{eq:R-prec}\\
& y\big(\etamin_t(s)+w\big)-x(s)\ge q. \label{eq:R-wait}
\end{align}
\end{subequations}
The envelopes are data at slot $t$, so every row of $\calR_t$ is a difference constraint. Let $\calR_{i,t}$ be the
system of plant $i$. A dispatch respects $\calR_{i,t}$ if at every slot $t$ it executes the slot-$t$ starts of an
integer point of $\calR_{i,t}$, with the round-robin assignment of units and the $k$-th tapped ladle in cast $k+q$. It
is robust zero-break if the schedule it executes satisfies \eqref{eq:zb} at the realized heat times, for all heat times
within the bounds.

At every slot $t$, Algorithm~1 of the paper builds every $\calR_{i,t}$ and solves the fleet LP over the $m_t$ blocks
$\calB_t$ that remain, with the variables $C$ and $C_\beta$, the rows
\begin{equation}\label{eq:rollcap}
\sum_{i\in\calI}\ell_{i,\beta}-\bar L_\beta\le C_\beta\le C,\qquad\beta\in\calB_t,
\end{equation}
and the plant rows $(x_i,y_i)\in\calR_{i,t}$. Its objective is $C+\frac1{m_t}\sum_{\beta\in\calB_t}C_\beta$ minus a small linear
reward for early starts after the window. It rounds an optimal vertex by the derandomized
rounding, with $\CLP$ replaced by the peak $C$ of the vertex, searches the combinations of shifted floors for a lower
planned peak, and executes the starts of slot $t$.

\begin{theorem}[Certified Rolling Dispatch]
(i) If $\calR_{i,0}$ is feasible and a dispatch respects $\calR_{i,t}$, then $\calR_{i,t}$ is feasible at every slot
$t$ and the dispatch is robust zero-break for plant $i$. (ii) If every $\calR_{i,0}$ is feasible, Algorithm~1 is robust
zero-break, and Theorem~3(i) holds at every slot with $m_t$ in place of $m$.
\end{theorem}

\subsection{The Order of Taps}

Heats on different furnaces may overtake one another, so the $k$-th tap need not come from the heat with the $k$-th
start. The envelopes still bound it.

\begin{lemma}\label{lem:taps}
Let $s_k$ be the $k$-th heat start. For all heat times within the bounds that agree with the observations up to slot
$t$, the $k$-th tap falls in $[\etamin_t(s_k),\,\bar\eta_t(s_k)]$.
\end{lemma}

\begin{proof}
A heat with start index $k'\le k$ starts at $s_{k'}\le s_k$ and taps by $s_{k'}+\dmax_t(s_{k'})\le\bar\eta_t(s_k)$,
so at least $k$ taps occur by $\bar\eta_t(s_k)$, and the $k$-th tap occurs by then. A heat with start index $k'\ge k$
starts at $s_{k'}\ge s_k$ and taps at or after $s_{k'}+\dmin_t(s_{k'})\ge\etamin_t(s_k)$, so at most $k-1$ taps occur
before $\etamin_t(s_k)$, and the $k$-th tap does not.
\end{proof}

\subsection{Proof of the Theorem}

\begin{proof}[Proof of Theorem 4]
\emph{Step 1: every integer point of $\calR_t$ is safe.} Let $p$ be an integer point of $\calR_t$, executed with the
round-robin assignment and first-in first-out ladles, and consider any heat times within the bounds that agree with the
observations up to slot $t$. Let $s_k$ be the $k$-th start of $p$.
\begin{itemize}[leftmargin=1.5em]
\item Furnaces. Row \eqref{eq:R-reuse} at $s=s_k$ gives $x(\bar\eta_t(s_k)+g-1)\le x(s_k-1)+\NE\le k-1+\NE$, so heat
$k+\NE$ starts at or after $\bar\eta_t(s_k)+g$. The heat started at $s_k$ taps by $s_k+\dmax_t(s_k)\le\bar\eta_t(s_k)$,
and heat $k+\NE$ is the next heat on its furnace, so \eqref{eq:zb-furn} holds.
\item Ladle precedence. Let cast $k+q$ start at $z$. Row \eqref{eq:R-prec} at $s=z$ gives
$x(\bar\theta_t(z))\ge y(z)-q\ge k$, so $s_k\le\bar\theta_t(z)$ and, as $\bar\eta_t$ is non-decreasing,
$\bar\eta_t(s_k)\le z$. By Lemma~\ref{lem:taps}, the $k$-th tap occurs by $z$, so the $k$-th tapped ladle, which cast
$k+q$ receives, is ready when the cast starts.
\item Waiting limit. Row \eqref{eq:R-wait} at $s=s_k$ gives $y(\etamin_t(s_k)+w)\ge x(s_k)+q\ge k+q$, so cast $k+q$
starts by $\etamin_t(s_k)+w$, which by Lemma~\ref{lem:taps} is at most $w$ slots after the $k$-th tap.
\item Casters. Rows \eqref{eq:F-busy} and \eqref{eq:F-nobreak} do not involve heat times, so Lemma~\ref{lem:order-stat}
and the round-robin assignment give \eqref{eq:zb-cast}.
\end{itemize}
Together with the rows of the waiting ladles, the schedule of $p$ thus satisfies \eqref{eq:zb} at these heat times.

\emph{Step 2: recursive feasibility.} Let $p_t$ be the integer point of $\calR_t$ whose slot-$t$ starts are executed.
New observations at slot $t+1$ shrink the residual intervals, so $\bar\eta_{t+1}\le\bar\eta_t$ and
$\etamin_{t+1}\ge\etamin_t$ pointwise, and hence $\bar\theta_{t+1}\ge\bar\theta_t$. Since the counters are
non-decreasing, each row of $\calR_{t+1}$ is implied at $p_t$ by the same row of $\calR_t$:
\[
x\big(\bar\eta_{t+1}(s)+g-1\big)\le x\big(\bar\eta_t(s)+g-1\big),\quad
x\big(\bar\theta_{t+1}(s)\big)\ge x\big(\bar\theta_t(s)\big),\quad
y\big(\etamin_{t+1}(s)+w\big)\ge y\big(\etamin_t(s)+w\big).
\]
The counters that $\calR_{t+1}$ fixes before $t+1$ are those of $p_t$, whose slot-$t$ starts were executed and whose
earlier starts were already fixed. Hence $p_t\in\calR_{t+1}$, and by induction from $\calR_{i,0}$, every
$\calR_{i,t}$ is feasible, so the dispatch is defined at every slot.

\emph{Step 3: conclusion of (i).} At the last slot $T$, the executed schedule coincides on the horizon with $p_T$, and
the realized heat times agree with all observations. By Step 1 applied at slot $T$, the executed schedule satisfies
\eqref{eq:zb} at the realized heat times. This holds for all heat times within the bounds, so the dispatch is robust
zero-break.

\emph{Step 4: part (ii).} The rows of $\calR_{i,t}$ are difference constraints with integer data, so by
\eqref{eq:floor} a common shift keeps all of them, and every combination of shifted floors that the search may choose is
an integer point of $\calR_{i,t}$. Algorithm~1 executes its starts with the round-robin assignment and first-in
first-out ladles, so it respects every $\calR_{i,t}$, and part (i) shows that it is robust zero-break. For the
guarantee, let $(C,(C_\beta),x,y)$ be an optimal vertex of the LP at slot $t$. Every row on the left of
\eqref{eq:rollcap} is tight, since a smaller $C_\beta$ would lower the objective, and let $\mu\le m_t$ rows on the
right be tight. The face count of Step 6 in Section~\ref{sec:thm3}, now with the $1+m_t$ variables $C$ and $C_\beta$,
gives $1+m_t+\sum_i\dim G_i\le m_t+\mu$, so at most $\mu-1\le m_t-1$ plants are fractional. Steps 1 to 5 and 7 of
the proof of Theorem~3(i) then hold for the $m_t$ remaining blocks, with $C$ in place of $\CLP$ and $C_\beta$ in place of
$\CLP_\beta$, and the search keeps the guarantee, since it accepts only a lower planned peak.
\end{proof}

\begin{thebibliography}{9}

\bibitem{bansal2008sicomp}
N.~Bansal, D.~Coppersmith, and M.~Sviridenko, ``Improved approximation algorithms for broadcast scheduling,''
\emph{SIAM J. Comput.}, vol.~38, no.~3, pp.~1157--1174, 2008.

\bibitem{durr2011algorithmica}
C.~D{\"u}rr and M.~Hurand, ``Finding total unimodularity in optimization problems solved by linear programs,''
\emph{Algorithmica}, vol.~59, no.~2, pp.~256--268, 2011.

\bibitem{mukhi2026tsg}
K.~Mukhi, G.~Loho, and A.~Abate, ``Exact characterization of aggregate flexibility via generalized polymatroids,''
\emph{IEEE Trans. Smart Grid}, vol.~17, no.~3, pp.~2131--2142, May 2026.

\bibitem{vujanic2016automatica}
R.~Vujanic, P.~Mohajerin Esfahani, P.~J. Goulart, S.~Mari{\'e}thoz, and M.~Morari, ``A decomposition method for large
scale MILPs, with performance guarantees and a power system application,'' \emph{Automatica}, vol.~67, pp.~144--156,
2016.

\bibitem{charikar2011soda}
M.~Charikar, A.~Newman, and A.~Nikolov, ``Tight hardness results for minimizing discrepancy,'' in \emph{Proc. ACM-SIAM
Symp. Discrete Algorithms}, 2011, pp.~1607--1614.

\end{thebibliography}

\end{document}
